\subsection{\highlight{Digital Models and Urban Twins} in Smart City Management}

The information support of smart city management processes is based on the development of an information model of the city and its processes \cite{x1}.
These models may vary in their degree of digital maturity and can be applied across different management horizons, as shown in Table \ref{table:maturity}.
At present, six stages can be distinguished based on the degree of integration of information and digital technologies into smart city management.

\begin{table}[h]
\centering
\small
\begin{tabularx}{\textwidth}{X | X | X | X}
\hline
& \textbf{Operational}& \textbf{Operations management}& \textbf{Strategic}\\
\hline
Form of management problem formulation& Task (SMART)& Project& Goals and target indicators\\
\cline{1-4}
Stage 1 “Informatization”
& \multirow{2}{=}{Industry-specific operational support systems (ERP)}
& \multirow{2}{=}{Project management systems (PMS)}
& Information and analytical support systems based on statistical reporting\\
\cline{1-1}\cline{4-4}
Stage 2 “Automation”
&
&
& BI systems, including automated generation of standard reports\\
\cline{1-4}
Stage 3 “Digitalization”& Systems incorporating city digital twins \cite{x2}& Digital project modeling systems (BIM, GIS)& Digital modeling of cities \cite{x5}\\
\cline{1-4}
Stage 4 “Intelligentization”& Digital twins with AI components \cite{x8}& Digital project models with AI components& Digital city modeling with predictive AI models \cite{x4,x6}\\
\cline{1-4}
Stage 5 “Autonomization”& Digital twins with humans completely excluded from the management of individual processes \cite{x12}& -& Digital models with AI-based adjustment of forecasts based on the current situation \cite{x10}\\
\cline{1-4}
Stage 6 “Generalization”& -& -& Automation of city goal-setting based on its role within the socio-economic system\\
\hline
\end{tabularx}
\caption{Levels of intelligent maturity in smart city management.}
\label{table:maturity}
\end{table}

The initial stages are \emph{informatization} and \emph{automation}, within which sector-specific information systems are developed.
These systems operate predominantly autonomously and support management tasks within individual sectors, such as transportation \cite{x3}, environmental management \cite{x7}, logistics, security \cite{x2}, and construction \cite{x4}.
At the early stages of informatization, the focus is primarily on operational and tactical management, owing to the greater degree of formalization of these processes.
Strategic planning is carried out primarily by experts based on fixed datasets.
A classic example is UrbanSim \cite{x5}, an urban development modeling system developed by Paul Waddell.
The platform is designed to analyze the interrelationships between land use, transportation processes, and environmental factors.
UrbanSim integrates models of population distribution, employment, development, changes in land value, and transportation accessibility, thereby enabling the dynamic forecasting of the spatial development of a territory \cite{x5}.
The system has been applied in modeling the Eugene–Springfield, Honolulu, and Salt Lake City metropolitan areas.

The \emph{digitalization} stage involves the creation of a digital representation of the object of management.
For sector-specific systems, the object of representation is the respective sectoral processes; at the operational level, it is projects; and for strategic planning, it is the city as a whole.
It should be noted that, at the operational level, the integration of individual systems gives rise to so-called city digital twins \cite{x2}, which are used to monitor the current operational situation.
At the strategic level, however, it is more appropriate to refer to digital city models, which support the development of long-term urban development plans.

The \emph{intelligentization} stage is characterized by the partial replacement of human decision-making by artificial intelligence (AI) models.
At present, such systems have the greatest impact in operational and, to a lesser extent, tactical activities, owing to the availability of the data required for model training \cite{x4}.
In design tasks, AI is being integrated into computer-aided design (CAD) systems \cite{x5}.
At the same time, the use of AI in digital city modeling for strategic management is primarily aimed at predicting the effects of proposed decisions \cite{x6,x10}.
One of the best-known examples of an integrated digital environment is the Virtual Singapore project, a national digital twin of Singapore \cite{x9}.
It is a highly detailed three-dimensional model of the city-state integrated with topographic and real-time operational data on buildings, infrastructure, population mobility, and environmental parameters.
The platform is used to model urban development scenarios, support transportation management, and prepare for emergency situations \cite{x9}.
A key feature of the project is the use of a unified collaborative environment for government agencies, research organizations, and businesses.

Problem formulation at all levels remains the least formalized component of the management process and consists primarily in defining the set of target indicators and their target values.

\emph{Autonomization} involves the exclusion of humans from the management of individual processes.
At the strategic planning level, this may manifest itself in adjusting the long-term forecast of the city's development based on the current situation \cite{x12}.

The final stage, \emph{generalization}, involves integrating cities with intelligent systems.
This is meaningful primarily in terms of better profiling the city's development goals based on information about the current state and development dynamics of other cities.

It can be seen that goal-setting remains the least formalized task.
It involves formulating goals as natural-language narratives and their native interpretation in the form of target indicators \cite{x13}.
At the same time, the long-term success of a city's development depends on the quality of its goal-setting, which necessitates the development of appropriate methods already at the early stages of \emph{digitalization}.